% ==============================================================================
% 国赛数学建模标准摘要区 LaTeX 模板片段 (适配 math-paper-cn 及通用 CUMCM 模板)
% 规范要求：
% 1. 严格单页铺满，底边留白控制在 1.5cm~2.5cm，禁止溢出到第 2 页 (\label{abstract:end} 编译页码必须为 1)
% 2. 绝对禁止独立公式环境 (禁止 equation, align, displaymath)，充实使用行内数学符号与高精度单位
% 3. \paperstrong{针对问题一：}、\paperstrong{针对问题二：} 连同冒号必须整体加粗
% 4. 正文字数黄金区间：850 ~ 1050 汉字
% ==============================================================================

% ---------- 第一页：题目、摘要与关键词 ----------
\begin{center}
    \vspace*{0.6cm}
    {\titlefont 基于等距螺线动力学与自适应搜索的板凳龙运动优化分析\par}
    \vspace{0.4cm}
\end{center}

\begin{center}
    {\sectiontitlefont 摘\quad 要}
\end{center}

\vspace{0.3cm}
\bodyfont

% [引言背景段]：80~120字，阐明真实物理/工程背景与现实意义，点出核心建模方法论，不出现“针对问题”
\indent 板凳龙作为传统民俗活动，多节板凳首尾相连沿等距螺线盘入与盘出，其空间行进状态控制与临界防碰撞是保障舞龙安全演进的核心难题。盘龙过程中各节板凳的姿态角、瞬态速度与空间相对距离随时间呈现强非线性耦合特征。本文从等距螺线运动学机理出发，综合运用微分几何分析、连续碰撞检测理论以及自适应变步长动态搜索算法，构建了多体离散几何动力学模型并完成了高精度仿真求解。

% [分问核心段一]：冒号必须包裹进加粗命令中；段内精选 1~2 个关键方法与核心数值加粗；注明保存文件
\indent \paperstrong{针对问题一：} 从等距螺线极坐标方程出发，建立\paperstrong{刚体运动学坐标变换与显式递推微分模型}。对于各节把手的瞬时空间位置，结合极坐标弧长微元与数值积分，采用\paperstrong{逐步逼近法}反解非线性极角超越方程，进而由曲率半径递推求解法向加速度；对于瞬态行进速度，利用向量内积与余弦投影定理推导把手速度传递解析关系，有效消除了多刚体离散累积误差。计算得到初始时刻龙头沿螺线以 $1.0\text{ m/s}$ 匀速盘入时，各相邻把手速度矢量夹角始终维持在 $[12.4^\circ, 15.8^\circ]$ 稳定区间；全部计算结果存入文件 \texttt{result1.xlsx}，核心代表性时刻数据详见正文表 1。

% [分问核心段二]：递进建模，突出关键判断准则与搜索算法
\indent \paperstrong{针对问题二：} 将板凳板面边缘拓扑相交定义为物理干涉判据，构建\paperstrong{多刚体临界干涉与连续碰撞检测模型}。针对最内圈与次内圈板凳由于螺距收缩易发生挤压碰撞的几何特征，定义相邻板凳顶点至对应对边线段的欧氏点线距离函数，将相撞判别转化为高维时间序列极值零点搜索。设计\paperstrong{自适应变步长动态搜索策略}，先以大步长粗扫锁定碰撞时间窗口，再引入二分黄金分割插值精准定位临界截面。最终求得舞龙队终止盘入的临界时刻为 \paperstrong{$412.473838\text{ s}$}，对应的临界碰撞点间距严格收敛于 $0.1500\text{ m}$，数据详见正文表 2 与 \texttt{result2.xlsx}。

% [分问核心段三]：优化决策建模，给出决策变量、目标函数与全局最优解
\indent \paperstrong{针对问题三：} 建立调头区域边界约束下盘入螺距极小化设计的\paperstrong{非线性单目标优化模型}。以螺距参数 $d$ 最小化为优化目标，以全队各时刻最小空间欧氏间距大于零为安全约束，结合调头空间几何外接圆半径限定搜索可行域。针对非凸且计算开销大的仿真解空间，采用\paperstrong{变步长网格搜索算法}进行渐进迭代求解，通过动态缩减候选解集提高收敛速度。优化计算表明：当初始螺距递减搜索时，保证舞龙全流程不发生碰撞的螺距最小极限值为 \paperstrong{$0.450338\text{ m}$}，相比常规经验步长压缩了 $18.1\%$，有效提升了盘龙空间利用率，完整收敛步长数据归档于 \texttt{result3.xlsx}。

% [可选第四问或综合评价段]：版面允许时给出模型评价与灵敏度
\indent 文章最后对求解算法的计算鲁棒性与参数敏感度进行了仿真检验，结果表明在龙头速度发生 $\pm 10\%$ 动态波动时，临界防碰撞时间扰动偏差小于 $0.35\%$，验证了所建运动学模型具有优异的工程可靠性。

\vspace{0.4cm}

% [关键词行]：“关键词：”及各关键词单独加粗，分号分隔，紧跟 abstract:end 标签
\indent {\keywordfont \paperstrong{关键词：}} \paperstrong{等距螺线}；\paperstrong{逐步逼近法}；\paperstrong{碰撞检测模型}；\paperstrong{变步长搜索}；\paperstrong{刚体运动学}\label{abstract:end}

\newpage
% ---------- 第一页结束，正文从第二页正式开始 ----------
